\section{曲线长度与内在距离}
\label{sec:riem-length-distance}

有了每个切空间中的长度，就可以把一条曲线各时刻的速率累计起来，得到曲线的长度。两点之间的距离则不是某一条曲线的属性，而要比较所有连接这两点的曲线。本节先定义长度并检查它与参数化无关，再用下确界定义内在距离，最后把它与环境空间中的直线距离比较。

\paragraph{曲线长度}
实际路线常有拐角，例如沿两段大圆弧拼成的航线。因此考虑\emph{分段光滑曲线}$\gamma:[a,b]\to M$：它连续，并且存在分点$a=t_0<t_1<\cdots<t_m=b$，使$\gamma$在每个$[t_{r-1},t_r]$上光滑。除有限个分点外，速度$\dot\gamma(t)\in T_{\gamma(t)}M$都有定义（第~\ref{chap:manifold-basics}章）。

\begin{definition}[曲线长度]
\label{def:riem-curve-length}
黎曼流形$(M,g)$上分段光滑曲线$\gamma:[a,b]\to M$的\term{曲线长度}（length of a curve）为
\begin{equation}
L_g(\gamma)=\int_a^b\lVert\dot\gamma(t)\rVert_g\,\mathrm dt
=\int_a^b\sqrt{g_{\gamma(t)}\bigl(\dot\gamma(t),\dot\gamma(t)\bigr)}\,\mathrm dt .
\label{eq:riem-curve-length}
\end{equation}
\end{definition}

若曲线落在一张坐标图内，写成坐标曲线$x(t)$，则被积函数是
\begin{equation*}
\sqrt{\sum_{i,j}g_{i,j}(x(t))\dot x^i(t)\dot x^j(t)}.
\end{equation*}
曲线跨越多张图时，紧区间$[a,b]$可以分成有限段，使每段落在一张图内；这一点来自$\gamma([a,b])$的紧性。式\eqref{eq:riem-metric-transform}保证在重叠处换图后被积函数不变，所以分段方式不影响结果。

\begin{example}[球面纬线的长度]
\label{ex:riem-latitude-length}
在例~\ref{ex:riem-sphere-metric}的坐标中取纬线$\gamma(t)=\Phi(\theta_0,t)$，$t\in[0,2\pi]$，严格说这条闭曲线有一点落在所取坐标域外，但单点不影响积分。坐标速度为$(\dot\theta,\dot\phi)=(0,1)$，于是$\lVert\dot\gamma\rVert_g^2=\sin^2\theta_0$，长度为$2\pi\sin\theta_0$。赤道$\theta_0=\pi/2$给出$2\pi$，靠近北极的纬线长度趋于零，与三维中圆的周长$2\pi\times(\text{半径}\sin\theta_0)$一致。
\end{example}

长度应当是路线本身的属性，不应取决于沿路线走得快还是慢。下面的结论说明定义满足这一要求。

\begin{theorem}[长度与参数化无关]
\label{thm:riem-length-reparametrization}
设$\gamma:[a,b]\to M$分段光滑，$\tau:[c,d]\to[a,b]$是分段光滑的单调满射，且$\tau'$除有限点外有定义并保持符号不变。则$L_g(\gamma\circ\tau)=L_g(\gamma)$。
\end{theorem}

\begin{proof}
由链式法则，$(\gamma\circ\tau)^{\cdot}(s)=\tau'(s)\dot\gamma(\tau(s))$，所以$\lVert(\gamma\circ\tau)^{\cdot}(s)\rVert_g=\lvert\tau'(s)\rvert\,\lVert\dot\gamma(\tau(s))\rVert_g$。若$\tau'\geq0$，作换元$t=\tau(s)$，$\mathrm dt=\tau'(s)\mathrm ds$，即得$\int_c^d\lVert\dot\gamma(\tau(s))\rVert_g\tau'(s)\,\mathrm ds=\int_a^b\lVert\dot\gamma(t)\rVert_g\,\mathrm dt$。若$\tau'\leq0$，换元后积分上下限互换，绝对值恰好抵消这一符号。分段光滑的情形逐段应用这一论证。
\end{proof}

定理要求$\tau$单调。若参数来回走，例如先前进再后退再前进，经过的路线被重复计入，长度会增加；这不是“换速度”，而是换了一条路线。

\paragraph{内在距离}
长度评价给定的一条路线，距离则要在所有允许路线中取最好的一条。这里“允许”指始终留在$M$中：穿过地球内部的直线不是地面上的路线。若两点之间所有路线的长度没有最小值，仍可以取下确界。

\begin{definition}[内在距离]
\label{def:riem-intrinsic-distance}
设$(M,g)$是黎曼流形，$p,q$属于$M$的同一个连通分支。定义$p,q$之间的\term{内在距离}（intrinsic distance）为
\begin{equation}
d_g(p,q)=\inf\bigl\{L_g(\gamma):\gamma:[a,b]\to M\text{分段光滑},\ \gamma(a)=p,\ \gamma(b)=q\bigr\}.
\label{eq:riem-intrinsic-distance}
\end{equation}
它也常称为测地距离。若$p,q$位于不同连通分支，约定$d_g(p,q)=+\infty$。
\end{definition}

这个集合非空，需要先有连接两点的分段光滑曲线。第~\ref{chap:manifold-basics}章证明连通流形是路径连通的，所用的路径在每个坐标球内由线段拼成，正是分段光滑曲线。因此在同一连通分支内，下确界是对非空集合取的，结果是有限非负数。

\begin{theorem}[内在距离是度量]
\label{thm:riem-distance-metric}
在黎曼流形$(M,g)$的每个连通分支上，$d_g$满足度量的三条公理：$d_g(p,q)\geq0$且仅当$p=q$时为零；$d_g(p,q)=d_g(q,p)$；$d_g(p,r)\leq d_g(p,q)+d_g(q,r)$。
\end{theorem}

\begin{proof}
对称性来自把曲线倒过来走：$t\mapsto\gamma(a+b-t)$连接$q$到$p$，由定理~\ref{thm:riem-length-reparametrization}长度不变。对三角不等式，任给$\varepsilon>0$，取连接$p,q$的曲线$\gamma_1$与连接$q,r$的曲线$\gamma_2$，使$L_g(\gamma_1)<d_g(p,q)+\varepsilon$，$L_g(\gamma_2)<d_g(q,r)+\varepsilon$。首尾相接得到连接$p,r$的分段光滑曲线，长度为两者之和，所以$d_g(p,r)<d_g(p,q)+d_g(q,r)+2\varepsilon$；令$\varepsilon\to0$即可。

剩下的是$p\neq q$时$d_g(p,q)>0$。取$p$附近的坐标图，使$p$的坐标为$x_p$，并取半径$\rho>0$的闭坐标球$\bar B=\{x:\lVert x-x_p\rVert_2\leq\rho\}$落在图内且不含$q$；这一步用到流形的Hausdorff性。$\bm G(x)$的最小特征值在紧集$\bar B$上连续且为正，因此有下界$\lambda>0$，使$\sum_{i,j}g_{i,j}(x)\xi^i\xi^j\geq\lambda\lVert\xi\rVert_2^2$对$x\in\bar B$成立。

任何从$p$到$q$的曲线都要离开$\bar B$，记它的坐标第一次到达球面$\lVert x-x_p\rVert_2=\rho$的时刻为$t^\ast$。在$[a,t^\ast]$上曲线留在$\bar B$中，所以
\[
L_g(\gamma)\geq\int_a^{t^\ast}\sqrt\lambda\,\lVert\dot x(t)\rVert_2\,\mathrm dt\geq\sqrt\lambda\,\bigl\lVert x(t^\ast)-x_p\bigr\rVert_2=\sqrt\lambda\,\rho .
\]
第二个不等式是“坐标曲线的欧氏长度不小于端点距离”。下界与曲线无关，所以$d_g(p,q)\geq\sqrt\lambda\,\rho>0$。
\end{proof}

于是每个连通分支$(M,d_g)$都是第~\ref{chap:metric-spaces}章意义下的度量空间，完备性、Cauchy列等概念都可以使用。若在闭坐标球上再取最大特征值的上界$\Lambda$，类似的估计给出$d_g$与坐标欧氏距离在局部互相控制，由此可以证明$d_g$诱导的拓扑恰为流形原有的拓扑。这里不给出完整证明；后文用到的只是这一结论：$d_g$意义下的收敛与流形拓扑中的收敛是同一回事。

\paragraph{内在距离与环境距离}
对嵌入$\mathbb R^D$并取诱导度量的流形，现在有两种距离：环境中的欧氏距离$\lVert q-p\rVert_2$，以及只允许沿$M$走的内在距离$d_g(p,q)$。二者总有确定的大小关系。对连接$p,q$的任一分段光滑曲线，把它看成$\mathbb R^D$中的曲线，由微积分基本定理与积分的三角不等式，
\[
\lVert q-p\rVert_2=\Bigl\lVert\int_a^b\dot\gamma(t)\,\mathrm dt\Bigr\rVert_2\leq\int_a^b\lVert\dot\gamma(t)\rVert_2\,\mathrm dt=L_g(\gamma),
\]
其中最后一步用的是诱导度量下$\lVert\dot\gamma\rVert_g=\lVert\dot\gamma\rVert_2$。对所有曲线取下确界，得到$\lVert q-p\rVert_2\leq d_g(p,q)$。

单位圆给出最简单的比较。圆心角为$\theta\in[0,\pi]$的两点之间，弦长是$2\sin(\theta/2)$，沿圆周的短弧长度是$\theta$，而任何连接这两点的圆上曲线都至少要扫过这段角度，所以$d_g=\theta$。$\theta=\pi/2$时两者分别约为$1.414$与$1.571$。展开$2\sin(\theta/2)=\theta-\theta^3/24+O(\theta^5)$可见，近处两种距离只差三阶小量，远处差别明显；对跖点之间弦长为$2$，内在距离为$\pi$。

\paragraph{下确界未必达到}
定义使用下确界而不是最小值，是因为最短路线可能不存在。在穿孔平面$M=\mathbb R^2\setminus\{0\}$上取欧氏度量，连接$p=(-1,0)$与$q=(1,0)$。沿横轴走到$(-\varepsilon,0)$，绕半径$\varepsilon$的上半圆到$(\varepsilon,0)$，再沿横轴到$q$，长度为
\[
(1-\varepsilon)+\pi\varepsilon+(1-\varepsilon)=2+(\pi-2)\varepsilon .
\]
令$\varepsilon\to0$，并结合上一段的下界$d_g(p,q)\geq\lVert q-p\rVert_2=2$，得到$d_g(p,q)=2$。长度恰为$2$的连接曲线只能是线段本身（可以改变速度），而线段经过被删去的原点，所以没有任何允许曲线达到下确界。

这个例子区分了两件事：两点之间存在连接曲线，只保证$d_g$有限；存在一条长度恰为$d_g$的曲线，是更强的要求。穿孔平面连通，第一件事成立，第二件事失败。最短曲线何时存在，要在第~\ref{sec:riem-shortest-paths}节借助完备性回答。
